\subsection{Sampled stationary branches}
Figure~\ref{fig:rotated-pentagon-branch-landscape} places this example
competitor in
the larger stationary-branch problem.  An independent floating-point producer
freezes \(3340\) raw two- and three-block words at the generic angle
\(\theta=\pi/20\), then solves the same words at \(73\) angles from \(0\) to
\(\pi/10\).  Its word count agrees with the exact certificate's count; the two
word sets were not compared independently.  This is an empirical explanation
of the finite classification, not another input to the proof.

Many raw words give numerically indistinguishable data on this grid.  To avoid
drawing the same curve on top of itself hundreds of times, the plot collapses
two words only when their statuses agree at all \(73\) angles and their sampled
actions agree after rounding to ten decimal places.  It then draws one complete
\(73\)-angle sequence for that group.  Thus every curve follows one word across
the whole grid; it is never assembled by selecting different words at different
angles.

\begin{figure}[p]
  \centering
  \includegraphics[width=0.98\linewidth]{working/rotated-regular-polygons/enumerated-kkt-branch-landscape.pdf}
  \caption{Sampled enumerated KKT-branch landscape.  Black marks the active
  branch and blue the example competitor \(\tau\), whose formula was calculated
  above.  Each gray curve is one of the
  other distinct sampled profiles defined in the preceding paragraph; the
  panel headings report the total number of distinct sampled profiles and the
  number of raw words they summarize.  A profile with no admissible sample
  produces no visible curve.  Only consecutive admissible samples are joined.
  Triangles just below the upper boundary mark curves continuing above the
  fixed cutoff \(6\); a gap means that the corresponding word was not
  numerically admissible at the missing samples.  The grid neither certifies
  feasibility sets nor proves the lower envelope.}
  \label{fig:rotated-pentagon-branch-landscape}
\end{figure}

\begin{table}[p]
\centering
\small
\textbf{Admissibility over the full 73-angle grid}\par\smallskip
\begin{tabular}{@{}lrrr@{}}
\toprule
sampled presence & two-block & three-block & all\\
\midrule
no admissible sample & \(795\) & \(1820\) & \(2615\)\\
admissible at every sample & \(50\) & \(155\) & \(205\)\\
one contiguous sampled run & \(105\) & \(414\) & \(519\)\\
multiple sampled runs & \(0\) & \(1\) & \(1\)\\
\bottomrule
\end{tabular}

\vspace{1.5em}
\textbf{Status at the right endpoint \(\theta=\pi/10\)}\par\smallskip
\begin{tabular}{@{}lrrr@{}}
\toprule
sampled status & two-block & three-block & all\\
\midrule
admissible & \(110\) & \(425\) & \(535\)\\
numerically inadmissible & \(770\) & \(1746\) & \(2516\)\\
indeterminate & \(70\) & \(219\) & \(289\)\\
solve failure & \(0\) & \(0\) & \(0\)\\
\bottomrule
\end{tabular}
\caption{Raw-word classification on the sampled grid.  The upper part
distinguishes words that are nowhere admissible on the grid, admissible
throughout, or admissible on one or more sampled runs; these classes are not
claims about feasibility-set topology.  The sole multiple-run word has
one numerically inadmissible grid point separating its two sampled admissible
runs.  The lower part records the producer's four-way status at
\(\theta=\pi/10\).  No solve failure occurred at any sampled angle.}
\label{tab:rotated-pentagon-sampled-presence}
\end{table}
\FloatBarrier
